% ======================================================================
%  Πρότυπο σειράς βοηθημάτων Μαθηματικών  (XeLaTeX)
%  Μεταγλώττιση:  xelatex main.tex  (δύο φορές για τα περιεχόμενα)
% ======================================================================

% ---------- Ρυθμίσεις που αλλάζουν συχνά ----------
\newcommand{\seriestitle}{Η Γλώσσα των Μαθηματικών}   % προσωρινός τίτλος σειράς (απόφαση 2 του business plan)
\newcommand{\booktitle}{Μαθηματικά ΣΤ' Δημοτικού}
\newcommand{\bookauthor}{Ιωάννης Ψωμιάδης}
\newif\ifdraft \drafttrue   % \draftfalse για την τελική έκδοση: κρύβει τις σημειώσεις συγγραφέα

% ---------- Σελίδα 17 x 24 cm ----------
\usepackage[paperwidth=170mm,paperheight=240mm,inner=20mm,outer=16mm,top=20mm,bottom=22mm,headsep=6mm,footskip=10mm]{geometry}

% ---------- Γλώσσα και γραμματοσειρές ----------
\usepackage{fontspec}
\usepackage{polyglossia}
\setdefaultlanguage{greek}
\setotherlanguage{english}
\defaultfontfeatures{Ligatures=TeX}
% Commissioner (Κώστας Μπαρτσόκας, άδεια OFL). Τα αρχεία είναι στον φάκελο fonts/.
\setmainfont{Commissioner}[Path=fonts/,Extension=.ttf,
  UprightFont=*-Regular,ItalicFont=*-Italic,BoldFont=*-Bold,BoldItalicFont=*-BoldItalic]
\setsansfont{Commissioner}[Path=fonts/,Extension=.ttf,
  UprightFont=*-Regular,ItalicFont=*-Italic,BoldFont=*-SemiBold,BoldItalicFont=*-BoldItalic]

% ---------- Μαθηματικά ----------
\usepackage{amsmath}
% Υποδιαστολή μέσα σε μαθηματικά: γράφουμε 3{,}5
\usepackage{unicode-math}      % σύμβολα από τη Latin Modern Math, ψηφία και γράμματα από την Commissioner
\setmathfont{Latin Modern Math}
\setmathfont{Commissioner-Regular.ttf}[Path=fonts/,range={up/num,up/latin,up/Latin,up/greek,up/Greek}]
\setmathfont{Commissioner-Italic.ttf}[Path=fonts/,range={it/latin,it/Latin,it/greek,it/Greek}]
\setmathfont{Latin Modern Math}[range={}]
% Κλάσματα μέσα στο κείμενο: σε μέγεθος που διαβάζει άνετα ένα παιδί (όχι τα μικροσκοπικά του \textstyle)
\AtBeginDocument{\let\oldfrac\frac
  \renewcommand{\frac}[2]{\mathchoice{\oldfrac{#1}{#2}}{\dfrac{#1}{#2}}{\oldfrac{#1}{#2}}{\oldfrac{#1}{#2}}}}
\newcommand{\tdot}{\mathord{.}}% τελεία χιλιάδων μέσα σε μαθηματικά, αν χρειαστεί

% ---------- Χρώματα: ένα spot χρώμα ----------
\usepackage{xcolor}
\definecolor{spot}{RGB}{22,84,150}       % το χρώμα της τάξης (ΣΤ': βαθύ μπλε· Ε': πορτοκαλί 184,78,24)
\definecolor{spotlight}{RGB}{228,238,249}
\definecolor{ink}{RGB}{40,40,40}
\definecolor{soft}{RGB}{242,242,242}
\definecolor{mid}{RGB}{110,110,110}

% ---------- Πακέτα σελιδοποίησης ----------
\usepackage{tcolorbox}
\tcbuselibrary{skins,breakable}
\usepackage{enumitem}
\usepackage{tabularx,booktabs,array,multicol}
\usepackage{needspace} % \Needspace{..} όπου ένας τίτλος πρέπει να μείνει με το πλαίσιο που ακολουθεί
\usepackage{xltabular} % πίνακες που σπάνε σε δύο σελίδες (επαναλαμβάνεται η επικεφαλίδα)
\setlength{\LTpre}{0.5em}\setlength{\LTpost}{0.3em}
% Ο πίνακας δεν χωρίζεται από τον τίτλο ή την πρόταση που προηγείται
\makeatletter\patchcmd{\LT@start}{\endgraf\penalty\z@\vskip\LTpre}{\endgraf\penalty\@M\vskip\LTpre}{}{\PackageWarning{preamble}{LT@start patch failed}}\makeatother
\usepackage{tikz}
\usetikzlibrary{arrows.meta,decorations.pathreplacing,calc}
\usepackage[nobottomtitles*]{titlesec} % κανένας τίτλος μόνος του στο κάτω μέρος της σελίδας
\renewcommand{\bottomtitlespace}{.14\textheight}
\usepackage{titletoc}
\usepackage{fancyhdr}
\usepackage[hidelinks]{hyperref}

\raggedbottom % οι σελίδες δεν «τεντώνονται» κάθετα για να γεμίσουν
\setlength{\parindent}{0pt}
\setlength{\extrarowheight}{2pt} % χώρος για κλάσματα μέσα σε πίνακες
\tolerance=2000
\emergencystretch=3em
\hbadness=10000
\setlength{\parskip}{0.55em}
\linespread{1.08}
\setlist{itemsep=0.25em,topsep=0.3em}
% Ποτέ μία μοναχική γραμμή παραγράφου στην αρχή ή στο τέλος σελίδας
\clubpenalty=10000 \widowpenalty=10000 \displaywidowpenalty=10000
\predisplaypenalty=10000
\makeatletter\@beginparpenalty=10000\makeatother % μια πρόταση που τελειώνει σε «:» μένει με τη λίστα/τον πίνακα που ακολουθεί

% ---------- Ρύθμιση σελιδοποίησης ανά κεφάλαιο ----------
% Το απόσταγμα μεταξύ παραγράφων ρυθμίζεται ανά κεφάλαιο (αρχείο selidopoiisi.tex),
% ώστε η τελευταία σελίδα κάθε κεφαλαίου να μη μένει με μία-δυο γραμμές.
\newlength{\basicparskip}\setlength{\basicparskip}{0.55em}
\newcommand{\tunepage}[1]{% #1 = κλειδί (π.χ. k13, e3)
  \ifcsname tune@#1\endcsname
    \setlength{\parskip}{\csname tune@#1\endcsname}%
  \else\setlength{\parskip}{\basicparskip}\fi}
\newcommand{\settune}[2]{\expandafter\def\csname tune@#1\endcsname{#2}}
\InputIfFileExists{selidopoiisi}{}{}

% ---------- Ενότητες (parts) και κεφάλαια ----------
\titleformat{\part}[display]{\centering\sffamily\color{spot}}{}{0pt}{\Huge\bfseries}
\titleformat{\chapter}[display]
  {\sffamily\bfseries\raggedright\hyphenpenalty=10000}
  {\color{spot}\large Κεφάλαιο \thechapter}
  {4pt}
  {\LARGE\textcolor{ink}}
\titlespacing*{\chapter}{0pt}{0pt}{10pt}
% Το χρώμα μπαίνει μέσα στον τίτλο (όχι στη μορφή): αλλιώς αφήνει «σημείο αλλαγής σελίδας» αμέσως μετά τον τίτλο
\titleformat{\section}{\sffamily\bfseries\large}{}{0pt}{\textcolor{spot}}
\titlespacing*{\section}{0pt}{1.2em}{0.4em}
\titleformat{\subsection}{\sffamily\bfseries}{}{0pt}{\textcolor{ink}}

\newcommand{\unitpart}[3]{%
  \clearpage\phantomsection\addcontentsline{toc}{part}{Ενότητα #1: #2}%
  \begingroup\renewcommand{\addcontentsline}[3]{}\part*{Ενότητα #1\\[6pt]\LARGE #2\\[10pt]\Large\normalfont\itshape «#3»}\endgroup
  \markboth{Ενότητα #1}{}}

\newcommand{\bookchapter}[3]{% {τίτλος}{αριθμός σχολικού κεφαλαίου}{τίτλος σχολικού}
  \setcounter{chapter}{\numexpr#2-1\relax}%
  \chapter{#1}\tunepage{k#2}%
  \setcounter{paradeigma}{0}\setcounter{askisi}{0}%
  \renewcommand{\theaskisi}{\thechapter.\arabic{askisi}}%
  {\sffamily\small\color{mid} Σχολικό βιβλίο: Κεφάλαιο #2, «#3»}\par\medskip}

\newcommand{\reviewchapter}[1]{%
  \chapter*{Επανάληψη ενότητας #1}\tunepage{e#1}%
  \addcontentsline{toc}{chapter}{\textit{Επανάληψη ενότητας #1}}%
  \markboth{Επανάληψη ενότητας #1}{Επανάληψη ενότητας #1}%
  \setcounter{askisi}{0}%
  \renewcommand{\theaskisi}{Ε#1.\arabic{askisi}}}

% Οι λευκές σελίδες πριν από μια δεξιά σελίδα μένουν εντελώς λευκές (χωρίς κεφαλίδα)
\makeatletter
\renewcommand{\cleardoublepage}{\clearpage\if@twoside\ifodd\c@page\else
  \hbox{}\thispagestyle{empty}\newpage\fi\fi}
\makeatother

% Περιεχόμενα
\setcounter{tocdepth}{1}
\makeatletter\renewcommand{\@pnumwidth}{2.2em}\makeatother
\contentsmargin{2.6em}
\titlecontents{part}[0pt]{\addvspace{11pt}\sffamily\bfseries\color{spot}}{}{}{\hfill\contentspage}[\addvspace{3pt}]
\titlecontents{chapter}[2.4em]{\addvspace{1.5pt}\sffamily}{\contentslabel{2.4em}}{}%
  {\titlerule*[0.7pc]{.}\contentspage}
\titlecontents{section}[2.4em]{\sffamily\small}{}{}{\titlerule*[0.7pc]{.}\contentspage}

% Κεφαλίδες
\pagestyle{fancy}
\fancyhf{}
\fancyhead[LE]{\sffamily\small\color{mid}\seriestitle}
\fancyhead[RO]{\sffamily\small\color{mid}\nouppercase{\leftmark}}
\fancyfoot[LE,RO]{\sffamily\small\thepage}
\renewcommand{\headrulewidth}{0pt}
\renewcommand{\chaptermark}[1]{\markboth{#1}{}}

% ---------- Τα μέρη κάθε κεφαλαίου ----------
\newcommand{\meros}[1]{\section*{#1}}

% 1. Η ανάγκη
\newtcolorbox{anagki}{enhanced,breakable,colback=spotlight,colframe=spotlight,
  borderline west={3pt}{0pt}{spot},arc=0pt,left=10pt,right=8pt,top=6pt,bottom=6pt,
  fontupper=\itshape}

% 2. Η λέξη
\newtcolorbox{lexi}[1]{enhanced,breakable,colback=white,colframe=spot,arc=2pt,boxrule=0.6pt,
  title={\sffamily\bfseries Η λέξη: #1},coltitle=white,colbacktitle=spot,left=8pt,right=8pt}

% 3. Η μετάφραση: πίνακας «στα ελληνικά / στα μαθηματικά»
\newenvironment{metafrasi}
  {\par\smallskip\renewcommand{\arraystretch}{1.45}\setlength{\extrarowheight}{5pt}%
   \setlength{\LTpre}{0pt}\setlength{\LTpost}{0pt}%
   \xltabular{\linewidth}{>{\raggedright\arraybackslash}X >{\centering\arraybackslash}p{0.44\linewidth}}
   \toprule
   \sffamily\bfseries Στα ελληνικά & \sffamily\bfseries Στα μαθηματικά\\ \midrule\endhead}
  {\bottomrule\endxltabular\par\smallskip}

% 4. Θεωρία: ορισμός, κανόνας, προσοχή
\newtcolorbox{orismos}[1][]{enhanced,breakable,colback=soft,colframe=ink,arc=2pt,boxrule=0.5pt,
  title={\sffamily\bfseries Ορισμός#1},coltitle=white,colbacktitle=ink,left=8pt,right=8pt}
\newtcolorbox{kanonas}[1][]{enhanced,breakable,colback=soft,colframe=ink,arc=2pt,boxrule=0.5pt,
  title={\sffamily\bfseries Κανόνας#1},coltitle=white,colbacktitle=ink,left=8pt,right=8pt}
\newtcolorbox{prosoxi}{enhanced,breakable,colback=white,colframe=spot,arc=2pt,boxrule=0.5pt,
  borderline west={3pt}{0pt}{spot},title={\sffamily\bfseries Προσοχή: συνηθισμένο λάθος},
  coltitle=spot,colbacktitle=white,left=8pt,right=8pt}

% 5. Λυμένα παραδείγματα: «τι ξέρω, τι θέλω, πώς τα συνδέω»
\newcounter{paradeigma}
\renewcommand{\theparadeigma}{\thechapter.\arabic{paradeigma}}
\newenvironment{paradeigma}[1][]
  {\refstepcounter{paradeigma}\par\medskip
   {\sffamily\bfseries\color{spot}Παράδειγμα \theparadeigma}\ifx&#1&\else{\sffamily\bfseries\ #1}\fi\par\nopagebreak}
  {\par\medskip}
\newcommand{\xero}[1]{\par\textbf{Τι ξέρω.} #1\par}
\newcommand{\thelo}[1]{\par\textbf{Τι θέλω.} #1\par}
\newcommand{\syndeo}[1]{\par\textbf{Πώς τα συνδέω.} #1\par}
\newcommand{\apantisi}[1]{\par\textbf{Απάντηση.} #1\par}
\newcommand{\elegxos}[1]{\par\textbf{Έλεγχος.} #1\par}

% 6. Ασκήσεις σε τρεις ζώνες
\newcommand{\zoni}[2]{\subsection*{\color{spot}#1\quad{\normalfont\small\color{mid}#2}}}
\newcommand{\zoniA}{\zoni{Το κατάλαβα}{εφαρμογή του κανόνα}}
\newcommand{\zoniB}{\zoni{Το χρησιμοποιώ}{προβλήματα από τη ζωή}}
\newcommand{\zoniC}{\zoni{Το σκέφτομαι}{ερωτήσεις που ανοίγουν δρόμο}}
\newcounter{askisi}
\renewcommand{\theaskisi}{\thechapter.\arabic{askisi}}
\newenvironment{askisi}
  {\refstepcounter{askisi}%
   \csname @beginparpenalty\endcsname=10000 % η εκφώνηση δεν χωρίζεται από τον πίνακα/σχήμα που ακολουθεί
   \begin{itemize}[leftmargin=2.9em,labelwidth=2.7em,labelsep=0.2em,align=left,itemsep=0pt,topsep=0.4em,partopsep=0pt]
   \item[\sffamily\bfseries\color{spot}\theaskisi]}
  {\end{itemize}}
\newcommand{\kouti}{\mathord{\setlength{\fboxsep}{1pt}\fbox{\rule{0pt}{1.5ex}\hspace{0.5em}}}} % κουτάκι για ψηφίο που λείπει
\newcommand{\kena}[1][2.5cm]{\underline{\hspace{#1}}}   % κενό για συμπλήρωση
\newcommand{\grammes}[1][2]{\par\foreach \i in {1,...,#1}{\vspace{1.4em}\noindent\textcolor{mid}{\rule{\linewidth}{0.3pt}}\par}}

% 7. Το παρέμβλημα του συγγραφέα
\newtcolorbox{paremvlima}{enhanced,breakable,colback=white,colframe=spot,arc=0pt,boxrule=0pt,
  borderline north={1pt}{0pt}{spot},borderline south={1pt}{0pt}{spot},
  left=4pt,right=4pt,fontupper=\small}

% 7β. Το «Από το εργοστάσιο» κάθε κεφαλαίου βρίσκεται χωριστά, στο ergostasio/fNN.tex, για εύκολη διόρθωση:
\newcommand{\ergostasio}[1]{\IfFileExists{ergostasio/f#1.tex}{%
  \meros{Από το εργοστάσιο}\begin{paremvlima}\input{ergostasio/f#1}\end{paremvlima}}{}}

% 8. Λύσεις (στο τέλος του βιβλίου)
\newenvironment{lyseiskef}[2]{\section*{Κεφάλαιο #1: #2}}{}
\newcommand{\lysi}[2]{\par\noindent{\sffamily\bfseries\color{spot}#1}\quad #2\par}

% ---------- Βοηθητικά ----------
\newcommand{\stub}[1]{\begin{tcolorbox}[colback=soft,colframe=soft,arc=2pt]\itshape\color{mid}#1\end{tcolorbox}}
\newcommand{\draftnote}[1]{\ifdraft{\par\sffamily\footnotesize\color{red}[Σημείωση συγγραφέα: #1]\par}\fi}
\newcommand{\euro}{€}
\newcommand{\symb}[1]{\textsf{#1}}

% ---------- Σχήματα κλασμάτων (Ενότητα 3) ----------
% \lorida[πλάτος cm]{ίσα μέρη}{χρωματισμένα}: λωρίδα χωρισμένη σε ίσα μέρη
\newcommand{\lorida}[3][5]{\begin{tikzpicture}[baseline=0.05cm]
  \pgfmathsetmacro{\lw}{#1/#2}%
  \foreach \i in {1,...,#2}{%
    \ifnum\i>#3\relax\def\lc{white}\else\def\lc{spot!35}\fi
    \draw[fill=\lc,draw=ink] ({(\i-1)*\lw},0) rectangle ({\i*\lw},0.42);}
\end{tikzpicture}}
% \pita[ακτίνα cm]{ίσα μέρη}{χρωματισμένα}: κύκλος (πίτα) χωρισμένος σε ίσα μέρη
\newcommand{\pita}[3][0.55]{\begin{tikzpicture}[baseline=-0.1cm]
  \foreach \i in {1,...,#2}{%
    \ifnum\i>#3\relax\else
      \fill[spot!35] (0,0) -- ({90-(\i-1)*360/#2}:#1) arc[start angle={90-(\i-1)*360/#2},end angle={90-\i*360/#2},radius=#1] -- cycle;\fi
    \ifnum#2>1\relax\draw[ink] (0,0) -- ({90-(\i-1)*360/#2}:#1);\fi}
  \draw[ink] (0,0) circle (#1);
\end{tikzpicture}}
% \kgrammi[cm ανά μονάδα]{μέγιστος ακέραιος}{παρονομαστής}{αριθμητής/ετικέτα,...}: αριθμογραμμή
\newcommand{\kgrammi}[4][4]{\begin{tikzpicture}[baseline=0cm]
  \draw[-{Latex},ink] (-0.2,0) -- ({#2*#1+0.4},0);
  \pgfmathtruncatemacro{\kn}{#2*#3}%
  \foreach \j in {0,...,\kn}{\draw[ink] ({\j*#1/#3},0.08) -- ({\j*#1/#3},-0.08);}
  \foreach \j in {0,...,#2}{\draw[ink,thick] ({\j*#1},0.15) -- ({\j*#1},-0.15) node[below]{\small\j};}
  \def\kmarks{#4}\ifx\kmarks\empty\else
  \foreach \m/\l in {#4}{\fill[spot] ({\m*#1/#3},0) circle (2pt) node[above=2pt]{\small\l};}\fi
\end{tikzpicture}}

% ---------- Διαγράμματα δεδομένων (Ενότητα 4) ----------
% \rabdogramma[cm ανά μονάδα]{ετικέτα/τιμή,...}{μέγιστο άξονα}{βήμα άξονα}
\newcommand{\rabdogramma}[4][0.1]{\begin{tikzpicture}[x=1cm,y=#1 cm]
  \foreach \y in {0,#4,...,#3}{\draw[mid,very thin] (0.3,\y) -- (0.25,\y) node[left,font=\scriptsize]{\y};}
  \foreach \l/\v [count=\i] in {#2}{\xdef\rlast{\i}%
    \fill[spot!55] (\i*1.45-0.38,0) rectangle (\i*1.45+0.38,\v);
    \node[above,font=\scriptsize] at (\i*1.45,\v) {\v};
    \node[below,font=\scriptsize,align=center,text width=1.4cm] at (\i*1.45,0) {\l};}
  \draw[ink] (0.3,#3) -- (0.3,0) -- (\rlast*1.45+0.7,0);
\end{tikzpicture}}
% \diplorabdo[cm ανά μονάδα]{ετικέτα/τιμή1/τιμή2,...}{μέγιστο}{βήμα}{υπόμνημα 1}{υπόμνημα 2}
\newcommand{\diplorabdo}[6][0.1]{\begin{tikzpicture}[x=1cm,y=#1 cm]
  \foreach \y in {0,#4,...,#3}{\draw[mid,very thin] (0.3,\y) -- (0.25,\y) node[left,font=\scriptsize]{\y};}
  \foreach \l/\v/\w [count=\i] in {#2}{\xdef\rlast{\i}%
    \fill[spot!70] (\i*1.7-0.5,0) rectangle (\i*1.7,\v);
    \fill[ink!35] (\i*1.7,0) rectangle (\i*1.7+0.5,\w);
    \node[below,font=\scriptsize,align=center,text width=1.6cm] at (\i*1.7,0) {\l};}
  \draw[ink] (0.3,#3) -- (0.3,0) -- (\rlast*1.7+0.8,0);
  \fill[spot!70] (\rlast*1.7+1.1,#3) rectangle ++(0.3,-#4/2); \node[right,font=\scriptsize] at (\rlast*1.7+1.4,#3-#4/4) {#5};
  \fill[ink!35] (\rlast*1.7+1.1,#3-#4) rectangle ++(0.3,-#4/2); \node[right,font=\scriptsize] at (\rlast*1.7+1.4,#3-#4*1.25) {#6};
\end{tikzpicture}}
% \eikona{πλήθος}: σειρά από εικονίδια (κυκλάκια) για εικονογράμματα
\newcommand{\eikona}[1]{\tikz[baseline=-0.6ex]{\foreach \i in {1,...,#1}{\fill[spot!70] (\i*0.36,0) circle (0.13);}}}
% \misoeikona: μισό εικονίδιο
\newcommand{\misoeikona}{\,\tikz[baseline=-0.6ex]{\fill[spot!70] (0,-0.13) arc[start angle=-90,end angle=90,radius=0.13] -- cycle;}}

% ---------- Δεκαδικοί (Ενότητα 5) ----------
% \plegma[πλευρά cm]{χρωματισμένα από 100}: τετράγωνο 10x10, χρωματίζονται στήλη-στήλη
\newcommand{\plegma}[2][2]{\begin{tikzpicture}[baseline=0.3cm,x=#1 cm/10,y=#1 cm/10]
  \pgfmathtruncatemacro{\pf}{#2/10}\pgfmathtruncatemacro{\pr}{#2-10*\pf}%
  \ifnum\pf>0 \fill[spot!45] (0,0) rectangle (\pf,10);\fi
  \ifnum\pr>0 \fill[spot!45] (\pf,10-\pr) rectangle (\pf+1,10);\fi
  \foreach \g in {1,...,9}{\draw[ink!40,very thin] (\g,0) -- (\g,10); \draw[ink!40,very thin] (0,\g) -- (10,\g);}
  \draw[ink] (0,0) rectangle (10,10);
\end{tikzpicture}}

% ---------- Ακέραιοι (Ενότητα 6) ----------
% \akgrammi[cm ανά μονάδα]{από}{έως}{θέση/ετικέτα,...}: αριθμογραμμή με αρνητικούς
\newcommand{\akgrammi}[4][0.55]{\begin{tikzpicture}[baseline=0cm,x=#1 cm]
  \draw[{Latex}-{Latex},ink] (#2-0.6,0) -- (#3+0.6,0);
  \foreach \j in {#2,...,#3}{\draw[ink] (\j,0.1) -- (\j,-0.1) node[below,font=\scriptsize]{$\j$};}
  \draw[ink,thick] (0,0.18) -- (0,-0.18);
  \def\kmarks{#4}\ifx\kmarks\empty\else
  \foreach \m/\l in {#4}{\fill[spot] (\m,0) circle (2.2pt) node[above=2pt,font=\small]{\l};}\fi
\end{tikzpicture}}
